\section*{NeurIPS Paper Checklist}

% DRAFT PASS -- items answered where the underlying information already
% exists in the project (design docs, protocol docs, project status notes)
% independent of whether the confirmatory model run has happened yet.
% Items that genuinely depend on results are left \answerTODO{} rather than
% guessed. Every drafted answer below should be re-checked once Sections
% 1, 2, 4-6 exist, since several justifications point at section numbers
% that aren't final yet.

\begin{enumerate}

\item {\bf Claims}
    \item[] Question: Do the main claims made in the abstract and introduction accurately reflect the paper's contributions and scope?
    \item[] Answer: \answerTODO{}
    \item[] Justification: \justificationTODO{} % Blocked: abstract/intro not yet drafted (depend on results).

\item {\bf Limitations}
    \item[] Question: Does the paper discuss the limitations of the work performed by the authors?
    \item[] Answer: \answerYes{}
    \item[] Justification: Section~\ref{sec:limitations} discusses the fixed intentionality level, fixed relationship context per family, the not-yet-piloted severity manipulation, the researcher-authored (not corpus-validated) obligation/severity taxonomy, the binary gender operationalization, and the single-API-aggregator model roster.

\item {\bf Theory assumptions and proofs}
    \item[] Question: For each theoretical result, does the paper provide the full set of assumptions and a complete (and correct) proof?
    \item[] Answer: \answerNA{}
    \item[] Justification: The paper presents an empirical vignette study; it contains no theorems or proofs.

\item {\bf Experimental result reproducibility}
    \item[] Question: Does the paper fully disclose all the information needed to reproduce the main experimental results of the paper?
    \item[] Answer: \answerTODO{}
    \item[] Justification: \justificationTODO{} % Blocked: the confirmatory run has not been executed, so we cannot yet confirm the disclosed design/prompt/schema (Sec.~\ref{sec:dataset}) is sufficient to reproduce results that don't exist yet. Revisit once Results is drafted -- the underlying design, prompt, and model roster are already fully specified, so this is likely answerable as Yes once the run happens, contingent on documenting any post-hoc analysis choices (e.g. hedge/refusal coding rule, categorical cutpoint) made during that run.

\item {\bf Open access to data and code}
    \item[] Question: Does the paper provide open access to the data and code, with sufficient instructions to faithfully reproduce the main experimental results?
    \item[] Answer: \answerYes{}
    \item[] Justification: The full vignette generation pipeline (design specification, generator script, generated 288-vignette set) and the data collection/validation scripts are maintained in a version-controlled repository. An anonymized copy will be linked in supplementary material at submission time to preserve double-blind review; the repository will be made public upon acceptance.

\item {\bf Experimental setting/details}
    \item[] Question: Does the paper specify all the training and test details necessary to understand the results?
    \item[] Answer: \answerYes{}
    \item[] Justification: Section~\ref{sec:dataset} and the (forthcoming) Measurement Protocol section specify the full factorial design, the fixed system prompt and JSON output schema, and the finalized 5-model/4-provider roster. Two calibration parameters -- the stability-pass repeat count $N$ and its sampling temperature -- are not yet finalized and must be pinned down and reported before the confirmatory run is treated as final; this is flagged explicitly rather than silently decided.

\item {\bf Experiment statistical significance}
    \item[] Question: Does the paper report error bars suitably and correctly defined or other appropriate information about the statistical significance of the experiments?
    \item[] Answer: \answerTODO{}
    \item[] Justification: \justificationTODO{} % Blocked: no experiment has been run yet, so no error bars or significance tests exist to report. The pre-registered analysis plan (Section~\ref{sec:planned-analysis}) specifies a mixed-effects model with scenario as a random effect and dispersion of fault_rating across stability-pass samples as the confidence measure -- fill in this answer once that analysis is actually run and reported in Section~\ref{sec:results}.

\item {\bf Experiments compute resources}
    \item[] Question: For each experiment, does the paper provide sufficient information on the computer resources needed to reproduce the experiments?
    \item[] Answer: \answerTODO{}
    \item[] Justification: \justificationTODO{} % Blocked: no run has been executed yet, so per-call/total API cost and wall-clock time are not yet known. All inference is via commercial/hosted APIs (OpenRouter), not self-hosted compute, so this section will report API call counts and cost rather than GPU/CPU-hours.

\item {\bf Code of ethics}
    \item[] Question: Does the research conducted in the paper conform, in every respect, with the NeurIPS Code of Ethics?
    \item[] Answer: \answerYes{}
    \item[] Justification: The study involves no human subjects or crowdsourced labor; all vignettes are synthetic constructions authored and internally reviewed by the research team against a documented writing-standards checklist (single-violation constraint excludes abuse, coercion, discrimination, or illegal-conduct content by design). All model access is via commercial APIs used within their published terms of service.

\item {\bf Broader impacts}
    \item[] Question: Does the paper discuss both potential positive societal impacts and negative societal impacts of the work performed?
    \item[] Answer: \answerYes{}
    \item[] Justification: Section~\ref{sec:impacts} discusses the positive impact (a reusable bias-detection methodology for LLM-as-judge systems deployed in relationally consequential contexts) and two negative/risk considerations (narrow patching of specific findings without addressing underlying representational bias; the binary gender operationalization's limits).

\item {\bf Safeguards}
    \item[] Question: Does the paper describe safeguards that have been put in place for responsible release of data or models that have a high risk for misuse?
    \item[] Answer: \answerNA{}
    \item[] Justification: The paper releases no model. The released dataset consists of synthetic, non-identifying relationship-conflict vignettes; per the writing-standards single-violation constraint, vignettes exclude abuse, coercion, threats, and illegal conduct by construction, so it poses no elevated misuse risk requiring a safeguard.

\item {\bf Licenses for existing assets}
    \item[] Question: Are the creators or original owners of assets used in the paper properly credited and are the license and terms of use explicitly mentioned and properly respected?
    \item[] Answer: \answerYes{}
    \item[] Justification: The study uses five third-party LLMs (Claude Sonnet, GPT-5-mini, Gemini 2.5 Flash, Llama 3.3 70B, DeepSeek V3.2) accessed via the OpenRouter inference API rather than any pre-existing scraped or repackaged dataset; each provider's API terms of service govern usage, no model weights are redistributed, and the originating paper/technical report for each model will be cited alongside its exact version identifier once the confirmatory run is executed.

\item {\bf New assets}
    \item[] Question: Are new assets introduced in the paper well documented and is the documentation provided alongside the assets?
    \item[] Answer: \answerYes{}
    \item[] Justification: The paper introduces a new 288-vignette factorial benchmark plus its generation pipeline, documented via a machine-readable design specification and three companion documents describing the schema, the post-drafting quality checklist, and human-readable renderings of every scenario. No human subjects or personal data are involved, so no consent question applies to the asset itself; license terms for the released dataset/code will be finalized before submission.

\item {\bf Crowdsourcing and research with human subjects}
    \item[] Question: For crowdsourcing experiments and research with human subjects, does the paper include the full text of instructions given to participants and screenshots, if applicable, as well as details about compensation?
    \item[] Answer: \answerNA{}
    \item[] Justification: The study involves no crowdsourced workers or external human subjects. An internal 20-vignette review sample was read through by members of the research team as part of standard content quality control, not as a human-subjects study.

\item {\bf Institutional review board (IRB) approvals or equivalent for research with human subjects}
    \item[] Question: Does the paper describe potential risks incurred by study participants, whether such risks were disclosed, and whether IRB approval was obtained?
    \item[] Answer: \answerNA{}
    \item[] Justification: The study does not involve human subjects research; the evaluated subjects are LLMs responding to synthetic vignettes, and internal content review was performed by the research team, not external participants.

\item {\bf Declaration of LLM usage}
    \item[] Question: Does the paper describe the usage of LLMs if it is an important, original, or non-standard component of the core methods in this research?
    \item[] Answer: \answerYes{}
    \item[] Justification: LLMs are the central object of study -- their fault-attribution judgments on the vignette set constitute the paper's primary dependent variable -- which is squarely a non-standard, load-bearing use requiring declaration, not an incidental writing aid. The (forthcoming) Measurement Protocol section fully specifies how each of the five evaluated models is prompted and how its output is parsed and scored. \textbf{Open item:} whether an LLM was also used to assist in drafting or tone-auditing vignette narrative text (referenced in project design docs) needs an explicit yes/no decision and, if yes, its own disclosure here, separate from the evaluation-subject usage above.

\end{enumerate}
